\documentclass[11pt]{article}
\usepackage{mathblatt}
% Präambel des Setzers (werkzeuge/setzer-praeambel.tex, wortgleich); dazu \gteil (Muster T2U)
\makeatletter
\setlength{\mbpfbe}{0mm}% keine Punkte (pruefung.md § 5: nur beim Prüfstein)
\newlength{\gEin}\setlength{\gEin}{0pt}
\let\gSchrift\relax
\renewcommand{\pfkreuzzeile}[1]{\par\noindent\hangindent3.2em\hangafter1\hspace*{1.6em}$\square$\ #1\par}
\newcommand{\gkreuz}[1]{$\square$\ #1\hspace{2em}}
\newcommand{\gkreuzzeile}[2]{\par\noindent\hangindent2.2em\hangafter1\mbox{#1\,$\square$}\ #2\par}% Buchstabe, Option: eine je Zeile (neben Skizze)
% links, was man ansieht, rechts, was man tut; Spaltengrenze überall an derselben Stelle (gemeinsam.md „Satz“)
% v0.17: Fußeintrag „Nr: Wert“ – kurz in einer Box (nie getrennt), lang umbrechbar mit der Nummer am Anfang
% (vorher \mbox: lange Lösungen liefen über den Rand, AS40 138 pt); robust, weil \fusshilfe in die .aux schreibt
\NewDocumentCommand{\gfuss}{m m}{\sbox0{#1:~#2}\ifdim\wd0<0.45\textwidth\mbox{#1:~#2}\else\mbox{#1:}\nobreak\ #2\fi}
\newsavebox{\gbox}\newlength{\gLinks}
% v0.10 (Lehrer 13.10.): keine starre Halbierung – die linke Spalte ist so breit wie die Figur (+ 6 mm), die
% Antwort steht direkt daneben; ist die Figur breiter als 0,62 Zeile, steht die Antwort darunter
\newcommand{\gzwei}[2]{\par\smallskip\noindent\sbox{\gbox}{#1}%
  \setlength{\gLinks}{\dimexpr\wd\gbox+6mm\relax}%
  \ifdim\gLinks>0.62\linewidth
    \usebox{\gbox}\par\smallskip\noindent\begin{minipage}{\linewidth}#2\end{minipage}%
  \else
    \begin{minipage}[t]{\gLinks}\vspace{0pt}\usebox{\gbox}\end{minipage}%
    \begin{minipage}[t]{\dimexpr\linewidth-\gLinks\relax}\vspace{0pt}\raggedright #2\end{minipage}%
  \fi\par}
% 16.10. (Lehrer, PA10 S. 1): Text links, Koordinatensystem rechts daneben; ist das System breiter als 0,58 Zeile,
% steht es wie bisher unter dem Text
% Passt das System knapp nicht mehr auf die Seite (bis 25 % zu hoch), wird es auf den Rest verkleinert statt
% die Seite halb leer zu lassen
\newlength{\gRest}\newlength{\gHoch}
\newcommand{\gtextks}[2]{\sbox{\gbox}{#1}%
  \setlength{\gHoch}{\dimexpr\ht\gbox+\dp\gbox\relax}%
  \ifdim\pagegoal<\maxdimen\setlength{\gRest}{\dimexpr\pagegoal-\pagetotal-4\baselineskip\relax}%
    \ifdim\gHoch>\gRest\ifdim\gRest>0.75\gHoch
      \sbox{\gbox}{\resizebox*{!}{\gRest}{\usebox{\gbox}}}\fi\fi\fi
  \ifdim\wd\gbox>0.58\linewidth #2\par\smallskip\noindent\gpassend{\usebox{\gbox}}\par
  \else\begin{minipage}[t]{\dimexpr\linewidth-\wd\gbox-5mm\relax}\raggedright #2\end{minipage}\hfill
    \begin{minipage}[t]{\wd\gbox}\vspace{0pt}\usebox{\gbox}\end{minipage}\par\fi}
% v0.9 (12.10., pruefung.md § 5 „Marke mit B, Bedingungen rechts“): drittes Argument von pfaufg = Bedingungen
% der Prüfung (v0.10: \pfohnetr = TR durchgestrichen, „FOR“), klein grau im rechten Rand auf der Grundlinie der Nummer, Zeilen mit \\
\newcommand{\gbedrand}[2]{\ifblank{#2}{}{\rlap{\hspace*{\dimexpr#1+2mm\relax}%
  \parbox[t]{15mm}{\raggedright\scriptsize\leavevmode\color{mbgrau}\strut#2}}}}
% Teil eines Bündels (a, b …): Marke links im Rand, Bedingungen rechts im Rand, auf der Zeile des Teils
\newcommand{\gteilrand}[2]{\llap{\mbpfrandmarke{#1}\hspace*{\mbpfnr}}\gbedrand{\linewidth}{#2}}
% Teil eines Bündels: Buchstabe in eigener Spalte, Inhalt eingerückt (\gEin = Einzug, damit die Spaltengrenze von
% \gzwei an derselben Stelle bleibt)
\newenvironment{gbteil}[3]{\par\smallskip\noindent\gteilrand{#2}{#3}\makebox[\mbpfteil][l]{\textbf{#1}}%
  \begin{minipage}[t]{\dimexpr\linewidth-\mbpfteil\relax}\setlength{\gEin}{\mbpfteil}\raggedright\strut\ignorespaces}%
 {\end{minipage}\par\addvspace{6pt}}% v0.18 (PG80 Nr. 1): Abstand, damit die Linie von a) nicht an b) rutscht
% Grafik höchstens so breit wie die Spalte (Bündel nebeneinander, Koordinatensystem ohne Karo)
\newsavebox{\gpbox}
\newcommand{\gpassend}[1]{\sbox{\gpbox}{#1}\ifdim\wd\gpbox>\linewidth\typeout{GPASSEND: verkleinert \the\wd\gpbox\space auf \the\linewidth}\resizebox{\linewidth}{!}{\usebox{\gpbox}}\else\usebox{\gpbox}\fi}
% Skizzenfeld (v0.10, Lehrer 13.10.): Karo ab dem linken Rand über die Zeile, „Skizze“ klein grau oben links im Feld
% (wie \pfskizzenfeld der Vorlage); der rechte Rand bleibt den Bedingungen
\newcommand{\gskizze}[1]{\par\smallskip\noindent\begin{tikzpicture}
  \clip (0,0) rectangle ({\linewidth-1mm},{#1*10mm});
  \draw[black!16,line width=0.3pt,step=5mm] (0,0) grid ({\linewidth},{#1*10mm});
  \node[anchor=north west,inner sep=2pt,font=\scriptsize,text=mbgrau,fill=white] at (0,{#1*10mm}) {Skizze};
  \end{tikzpicture}\par}
\RenewDocumentEnvironment{pfaufg}{m m m}{%
  \begin{lrbox}{\mbaufgabebox}\begin{minipage}[t]{\linewidth}%
  \gSchrift\noindent\hspace*{\gEin}\mbpfrandmarke{#1}\gbedrand{\linewidth-\gEin}{#3}%
  \begin{minipage}[t]{\mbpfnr}\strut\textbf{#2}\end{minipage}%
  \begin{minipage}[t]{\dimexpr\linewidth-\gEin-\mbpfnr\relax}\raggedright\strut\ignorespaces}%
 {\end{minipage}%
  \end{minipage}\end{lrbox}%
  \setlength{\mbaufgabehoehe}{\dimexpr\ht\mbaufgabebox+\dp\mbaufgabebox\relax}%
  \ifdefempty{\gkopfnach}{\par\addvspace{5pt}\Needspace*{\mbaufgabehoehe}}{%
    % Stufenkopf, der mit dieser (ersten) Aufgabe zusammenbleibt: Platz für beide verlangen, dann den Kopf setzen
    % (12.10., Quadratische: ein Kopf stand allein unten, die Aufgabe mit Koordinatensystem auf der nächsten Seite)
    % v0.17 (AS40 „Terme deuten“ allein unten): hinter dem Kopf steht das Sprungziel (\special, ein Knoten, nach
    % dem die Seite brechen durfte) – darum \nobreak nach dem Kopf; Vorhalt 56 pt = Kopf mit Linie und Abständen
    \par\Needspace*{\dimexpr\mbaufgabehoehe+56pt\relax}\gkopfnach\par\nobreak\gdef\gkopfnach{}\addvspace{5pt}}%
  \noindent\usebox{\mbaufgabebox}\par\addvspace{5pt}}
% v0.18 (plan.md § 10, 15.10. l): Abschnittsbuchstabe A, B, C vor dem Namen im Kopf – als Vorsatz, damit der
% Name im TeX gleich bleibt (Bestellkarte, Sprungziele lesen ihn)
\def\gbstakt{}\newcommand{\gbstvor}[1]{\gdef\gbstakt{#1}}
\NewCommandCopy\gpfstufekopfalt\pfstufekopf
\RenewDocumentCommand{\pfstufekopf}{o m m}{\IfValueTF{#1}{\gpfstufekopfalt[#1]}{\gpfstufekopfalt}%
  {\ifx\gbstakt\empty\else\gbstakt\enspace\fi #2}{#3}\gdef\gbstakt{}}
% v0.18 (16.10. e): Begriff grau, ohne Nummer, ohne Rahmen über der ersten Aufgabe der Stufe
\newcommand{\gbegriff}[1]{\par\nobreak\noindent\hspace*{\mbpfnr}\parbox[t]{\dimexpr\linewidth-\mbpfnr\relax}{\raggedright\small\color{mbgrau}#1}\par\nobreak\addvspace{3pt}}
% v0.18 (16.10. g): Lösungsblatt – Kopf mit kräftigem Strich, Abschnitt klein grau als Zwischenzeile
\newcommand{\glkopf}[1]{\noindent{\large\bfseries Lösungen \textperiodcentered{} #1}\par\vspace{3pt}\hrule height 1.2pt\par\medskip}
\newcommand{\glzwischen}[1]{\par\addvspace{8pt}\Needspace*{4\baselineskip}\noindent{\footnotesize\color{mbgrau}#1}\par\nobreak\vspace{1pt}}
% v0.18 (15.10. o, w): Blattende – Kern ohne Rahmen, dann graue Zeile „Mehr: n Aufgaben“
\newcommand{\gkern}[1]{\par\addvspace{14pt}\Needspace*{6\baselineskip}\noindent{\bfseries Kern}\par\nobreak\smallskip
  {\small\raggedright #1\par}}
\newcommand{\gmehr}[1]{\par\addvspace{10pt}\noindent{\footnotesize\color{mbgrau}Mehr: #1 Aufgaben}\par}
% Stufenkopf vormerken: \gkopfvor{<Kopf>} – die nächste pfaufg setzt ihn, sobald beide zusammen Platz haben
\def\gkopfnach{}
\newcommand{\gkopfvor}[1]{\gdef\gkopfnach{#1}}
% Teilaufgabe mit Skizze links in gzwei: Buchstabe in eigener Spalte, darunter Text und Skizze
\newcommand{\gteil}[2]{\makebox[\mbpfteil][l]{\textbf{#1}}\parbox[t]{\dimexpr0.5\textwidth-\gEin-\mbpfnr-5mm-\mbpfteil\relax}{\raggedright #2}}
% Stufenkopf der Vorlage (mathblatt.sty 2026-10-12: Linie, Abstand, 13 pt, \pfjahre) ohne ihren Vorhalt von
% 14 Zeilen: der Setzer fordert vor dem Kopf 0,16 Seitenhöhe (Kopf + erste Aufgabe), mehr nicht
\setlength{\mbpfstufevorhalt}{0pt}
\makeatother
\newcommand{\karo}[1]{\par\smallskip\noindent\begin{tikzpicture}
  \clip (0,0) rectangle ({\linewidth-1mm},{#1*10mm});
  \draw[black!16,line width=0.3pt,step=5mm] (0,0) grid ({\linewidth},{#1*10mm});
  \end{tikzpicture}\par}
\newcommand{\kk}[1]{$\square$\,#1\hspace{0.9em}}
\newcommand{\linie}[1]{\,{\color{black!45}\rule[-1pt]{#1}{0.4pt}}\,}
\newcommand{\pkt}{\textperiodcentered{}}
% v0.14 (14.10.): Stufenkopf mit fünf Punkten (Option = Jahre für die Bestellkarte) oder „zuletzt 19“
\NewDocumentCommand{\gpunkte}{o m}{#2}
\newcommand{\gzuletzt}[1]{zuletzt #1}
% v0.7 (12.10.): Schülerfrage unter dem Stufenkopf; Ankreuzoption als Bild (Buchstabe und Kästchen darunter)
\newcommand{\gfrage}[1]{\par\nobreak\noindent\hspace*{\mbpfnr}{\small\itshape #1}\par\nobreak\addvspace{2pt}}
\newlength{\gbildbreite}\setlength{\gbildbreite}{0.32\linewidth}% Setzer setzt sie je Reihe (vier Bilder: 0,24)
\newcommand{\gbildoption}[3]{\begin{minipage}[t][#3][b]{\gbildbreite}\centering #2\par\smallskip\mbox{#1\,$\square$}\end{minipage}}
% Lineare (12.10.): Bild mit Buchstabe ohne Kästchen (Diagramm zuordnen); Aussage mit „wahr / falsch“ zum Ankreuzen
\newcommand{\gbildname}[3]{\begin{minipage}[t][#3][b]{\gbildbreite}\centering #2\par\smallskip\mbox{\textbf{#1}}\end{minipage}}
\newcommand{\gwahrfalsch}[2]{\par\noindent\hangindent2.2em\hangafter1\makebox[2.2em][l]{(#1)}#2\hfill\mbox{wahr $\square$\quad falsch $\square$}\par}

\newcommand{\gunter}[2]{\par\addvspace{6pt}\noindent{\bfseries #1\enspace #2}\par\nobreak\vspace{2pt}}
\newcommand{\gsternzeile}{\par\vspace{-4pt}\noindent{\footnotesize\color{mbgrau}* nur FOR}\par\medskip}
% 17.10.: Skizze oben bündig mit der ersten Textzeile (beide Spalten hängen an der Oberkante); der Rechenraum (Karo)
% beginnt direkt unter dem Text, links neben dem Graph, und läuft unter dem Graph über die ganze Zeile weiter
% (Streifen links = auf 5 mm aufgerundete Höhe des Graphen unter dem Text, Raster durchgehend; die Fläche des Streifens
% wird von der Karohöhe unter dem Graph abgezogen, auf 0,5 cm aufgerundet – die Gesamtfläche Rechenraum bleibt mindestens #3 × Zeile); #3 = Zeilen Karo in cm
% unter dem Graph (0 = kein Karo). Kein Verkleinern des Graphen.
\newsavebox{\gtbox}\newlength{\gWl}\newlength{\gDh}
\newcommand{\gkaroohne}[1]{\nointerlineskip\noindent\begin{tikzpicture}
  \clip (0,0) rectangle ({\linewidth-1mm},{#1*10mm});
  \draw[black!16,line width=0.3pt,step=5mm] (0,0) grid ({\linewidth},{#1*10mm});
  \end{tikzpicture}\par}
\newcommand{\gtextkaro}[3]{\sbox{\gbox}{#1}%
  \ifdim\wd\gbox>0.58\linewidth #2\par\smallskip\noindent\gpassend{\usebox{\gbox}}\par\ifnum#3>0 \karo{#3}\fi
  \else
    \setlength{\gWl}{\dimexpr\linewidth-\wd\gbox-5mm\relax}%
    \sbox{\gtbox}{\begin{minipage}[t]{\gWl}\raggedright #2\par\end{minipage}}%
    \setlength{\gDh}{\dimexpr\ht\gbox+\dp\gbox-\ht\gtbox-\dp\gtbox-6pt\relax}%
    \leavevmode
    \ifnum#3>0\ifdim\gDh>0pt
      \pgfmathsetmacro{\gxs}{floor(\gWl/5mm)*5mm}\pgfmathsetmacro{\gys}{ceil(\gDh/5mm)*5mm}%
      \begin{minipage}[t]{\gWl}\usebox{\gtbox}\par\vspace{6pt}\nointerlineskip\noindent
        \begin{tikzpicture}
        \clip (0,0) rectangle (\gxs pt,\gys pt);
        \draw[black!16,line width=0.3pt,step=5mm] (0,0) grid (\gxs pt,\gys pt);
        \end{tikzpicture}\par\end{minipage}%
    \else\usebox{\gtbox}\fi\else\usebox{\gtbox}\fi
    \hfill\begin{minipage}[t]{\wd\gbox}\vspace{-\dimexpr\ht\strutbox+2.4ex\relax}\usebox{\gbox}\end{minipage}\par
    \ifnum#3>0 \ifdim\gDh>0pt
      \pgfmathsetmacro{\gnn}{max(0,ceil(2*(#3-(\gys/28.4528)*(\gxs/\linewidth)))/2)}%
      \ifdim\gnn pt>0pt\gkaroohne{\gnn}\fi
    \else\karo{#3}\fi\fi
  \fi}
\begin{document}

\pfheftstil
\pflaufkopf[P10 \pkt{} Prüfung]{Fläche und Umfang}{PG40}\pfseitevon
\pfheftkopf{Fläche und Umfang}{}
\pfstufenzeile{Umfang und Flächeninhalt von Figuren \pkt{} Flächeninhalt von Dreiecken \pkt{} Flächeninhalt von Trapezen \pkt{} Flächeninhalt von Drachen \pkt{} Kreis}
\gsternzeile

% ===================== A Umfang und Flächeninhalt von Figuren =====================
\gkopfvor{\gbstvor{A}\pfstufekopf{Umfang und Flächeninhalt von Figuren}{}}
% DRUCK 2014-OS-K5d deckt 2014-OS-K5d
\begin{pfaufg}{nach~P10\\’14}{1.}{}
\gtextkaro{\begin{tikzpicture}[scale=0.7,line width=0.7pt,font=\small]\fill[black!10] (0,4.4) -- (2.2,4.4) -- (2.2,1.1) -- (1.9,0) -- (0.3,0) -- (0,1.1) -- cycle;\fill[black!22] (0.3,3.6) rectangle (1.9,1.1);\draw (0,4.4) -- (2.2,4.4) -- (1.9,3.6) -- (0.3,3.6) -- cycle;\draw (0,4.4) -- (0,1.1) (2.2,4.4) -- (2.2,1.1);\draw (0.3,3.6) -- (0.3,0) (1.9,3.6) -- (1.9,0);\draw[dashed,line width=0.5pt] (0,1.1) -- (2.2,1.1);\draw (0,1.1) -- (0.3,0) -- (1.9,0) -- (2.2,1.1);\draw[line width=0.5pt] (1.1,4.4) -- (1.1,3.6);\draw[line width=0.4pt] (1.1,3.75) -- (1.25,3.75) -- (1.25,3.6);\draw[line width=0.5pt] (1.1,1.1) -- (1.1,0) node[pos=0.4,left,inner sep=1pt] {$h$};\draw[line width=0.4pt] (1.1,0.15) -- (1.25,0.15) -- (1.25,0);\draw[line width=0.4pt,<->] (0,4.65) -- (2.2,4.65) node[midway,above,inner sep=1pt] {110};\draw[line width=0.4pt,<->] (2.45,4.4) -- (2.45,1.1) node[midway,right,inner sep=1pt] {165};\draw[line width=0.4pt,<->] (2.42,1.02) -- (2.12,-0.08) node[midway,right,inner sep=2pt] {57};\draw[line width=0.4pt,<->] (0.3,-0.25) -- (1.9,-0.25) node[midway,below,inner sep=1pt] {80};\node[font=\scriptsize,color=mbgrau,anchor=north] at (1.1,-0.75) {(Maße in cm; nicht maßstabsgerecht)};\end{tikzpicture}}{Bäckermeister Neumann lässt eine Kühlvitrine in Form eines Prismas bauen (Maße in cm). Die Rückwand der Vitrine soll aus einer rechteckigen Platte ausgesägt werden. Es steht eine 1,8 m² große Platte zur Verfügung. Sie ist 1,20 m breit. Kann die Platte verwendet werden? Entscheide mit Hilfe einer Rechnung.}{4}
\fusshilfe{\gfuss{1}{nein}}
\end{pfaufg}
% DRUCK 2017-OS-K3a deckt 2017-OS-K3a
% DRUCK 2017-OS-K3b deckt 2017-OS-K3b
\begin{pfaufg}{nach~P10\\’17}{2.}{}
\gtextkaro{\begin{tikzpicture}[line width=0.8pt,font=\small]\fill[black!8] (0,0) -- (3.6,0) arc (-90:90:0.9) -- (0,1.8) arc (90:270:0.9) -- cycle;\draw (0,0) -- (3.6,0) arc (-90:90:0.9) -- (0,1.8) arc (90:270:0.9) -- cycle;\draw[dashed,line width=0.4pt] (0,0) -- (0,1.8) (3.6,0) -- (3.6,1.8);\draw[line width=0.4pt,<->] (-1.25,0) -- (-1.25,1.8) node[midway,left,inner sep=2pt] {$a = 5$ m};\draw[line width=0.4pt,<->] (0,-0.4) -- (3.6,-0.4) node[midway,below,inner sep=2pt] {$b = 10$ m};\draw[line width=0.4pt] (0,-0.25) -- (0,-0.55) (3.6,-0.25) -- (3.6,-0.55);\node[font=\scriptsize,color=mbgrau] at (1.8,-1.25) {(nicht maßstabsgerecht)};\end{tikzpicture}}{Familie Sommer hat einen Swimmingpool gebaut. Er hat folgende Grundfläche:\begin{gbteil}{a)}{}{}Gib an, aus welchen Teilflächen sich die Grundfläche des Swimmingpools zusammensetzt.\end{gbteil}}{3}
\begin{gbteil}{b)}{}{}
Berechne die Größe der Grundfläche des Swimmingpools.
\karo{5}
\par\medskip\noindent\hfill A $=$\linie{25mm}
\end{gbteil}
\fusshilfe{\gfuss{2}{a) Rechteck und zwei Halbkreise; b) 69,6 m²}}
\end{pfaufg}

% ===================== B Flächeninhalt von Dreiecken =====================
\gkopfvor{\gbstvor{B}\pfstufekopf{Flächeninhalt von Dreiecken}{}}
% DRUCK 2020-OS-K7b deckt 2020-OS-K7b
\begin{pfaufg}{nach~P10\\’20}{3.}{}
\gtextkaro{\begin{tikzpicture}[line width=0.7pt,font=\small]\coordinate (A) at (0,0);\coordinate (B) at (1.3,3.6);\coordinate (C) at (0,3.6);\coordinate (Q) at (0,2.01);\draw (A) -- (B) -- (C) -- cycle;\draw[densely dashed,line width=0.5pt] (Q) -- (B);\fill (Q) circle (1.2pt);\pfrw[0.3]{(C)}{-90}\node[anchor=north] at (A) {$A$};\node[anchor=south west] at (B) {$B$};\node[anchor=south east] at (C) {$C$};\node[anchor=east] at (Q) {$Q$};\node[font=\scriptsize,color=mbgrau,anchor=north] at (current bounding box.south) {(nicht maßstabsgerecht)};\end{tikzpicture}}{Gegeben ist das rechtwinklige Dreieck $ABC$. Die Strecke $\overline{BC}$ ist 12~m lang und die Strecke $\overline{AC}$ hat eine Länge von 34~m. Der Punkt $Q$ liegt auf der Strecke $\overline{AC}$. Der Flächeninhalt des Dreiecks $QBC$ beträgt 90~m². Berechne die Länge der Strecke $\overline{QC}$.}{4}
\par\medskip\noindent\hfill QC $=$\linie{25mm}
\fusshilfe{\gfuss{3}{15 m}}
\end{pfaufg}\vspace{-6pt}
% DRUCK 2015-OS-K6c deckt 2015-OS-K6c
\begin{pfaufg}{nach~P10\\’15}{4.}{}
\gtextkaro{\begin{tikzpicture}[x={(1.25cm,0cm)},y={(0.62cm,0.42cm)},z={(0cm,1cm)},line width=0.7pt]\draw (0,0,0) -- (3,0,0) -- (3,3,0);\draw (0,0,0) -- (1.5,1.5,2.7) -- (3,0,0) (3,3,0) -- (1.5,1.5,2.7);\draw[dashed,line width=0.5pt] (0,0,0) -- (0,3,0) -- (3,3,0) (0,3,0) -- (1.5,1.5,2.7);\node[font=\scriptsize,color=mbgrau,anchor=north] at (current bounding box.south) {(nicht maßstabsgerecht)};\end{tikzpicture}}{Tim hat ein Modell der Karlsruher Pyramide gebaut, einer quadratischen Pyramide. Sein Modell ist 0,68 m hoch. Die Grundkante ist im Modell 0,80 m lang. Die Seitenflächen der Modellpyramide baut Tim aus dünnem Sperrholz. Weise nach, dass er für eine Seitenfläche ca. 0,32 m² Sperrholz braucht.}{6}
\fusshilfe{\gfuss{4}{0,32 m²}}
\end{pfaufg}\vspace{-6pt}
% DRUCK 2022-OS-K5e deckt 2022-OS-K5e
\begin{pfaufg}{nach~P10\\’22}{5.\textsuperscript{*}}{}
\gtextkaro{\begin{tikzpicture}[baseline=(current bounding box.north),line width=0.7pt,font=\small]\draw (0,0) -- (7.11,0) -- (4.8,2.96) -- cycle;\draw[dashed] (4.8,2.96) -- (4.8,0);\pfrw[0.25]{(4.8,0)}{90}\node[anchor=north east,inner sep=1.5pt] at (0,0) {$A$};\node[anchor=north west,inner sep=1.5pt] at (7.11,0) {$B$};\node[anchor=south,inner sep=1.5pt] at (4.8,2.96) {$C$};\node[anchor=north,inner sep=1.5pt] at (4.8,0) {$D$};\node[anchor=west,inner sep=1.5pt] at (4.8,1.3) {$h_c$};\pfwinkel[0.8][0.62]{(7.11,0)}{128}{180}{52^\circ}\end{tikzpicture}}{Gegeben ist ein Dreieck $ABC$. Das Dreieck ist nicht rechtwinklig. Die Höhe $h_c$ beträgt ca. 7,4~m, die Strecke $\overline{AD}$ ist ca. 12,0~m lang. Ermittle den Flächeninhalt des Dreiecks $ABC$.}{5}
\par\medskip\noindent\hfill A $=$\linie{25mm}
\fusshilfe{\gfuss{5}{65,8 m²}}
\end{pfaufg}\vspace{-6pt}

% ===================== C Flächeninhalt von Trapezen =====================
\gkopfvor{\gbstvor{C}\pfstufekopf{Flächeninhalt von Trapezen}{}}
% DRUCK 2023-OS-K2b deckt 2023-OS-K2b
\begin{pfaufg}{nach~P10\\’23}{6.}{}
\gtextkaro{\begin{tikzpicture}[line width=0.7pt,font=\small]\draw (0,0) -- (5.16,0) -- (4.08,1.6) -- (1.08,1.6) -- cycle;\draw[densely dotted,line width=0.6pt] (1.08,1.6) -- (1.08,0);\pfrw[0.2]{(1.08,0)}{90}\node[anchor=west,font=\footnotesize] at (1.12,0.85) {$8{,}00$ m};\pfwinkel[0.55][0.6]{(0,0)}{0}{56}{}\node[font=\scriptsize] at (0.74,0.4) {$56^\circ$};\pfwinkel[0.55][0.6]{(5.16,0)}{124}{180}{}\node[font=\scriptsize] at (4.42,0.4) {$56^\circ$};\draw[line width=0.4pt,|-|] (1.08,1.95) -- (4.08,1.95) node[midway,above,inner sep=1pt,font=\footnotesize] {$15{,}00$ m};\draw[line width=0.4pt,|-|] (0,-0.35) -- (5.16,-0.35) node[midway,below,inner sep=1pt,font=\footnotesize] {$25{,}80$ m};\node[font=\scriptsize,color=mbgrau,anchor=north] at (current bounding box.south) {(nicht maßstabsgerecht)};\end{tikzpicture}}{Gegeben ist das folgende Trapez. Berechne den Flächeninhalt des Trapezes.}{4}
\par\medskip\noindent\hfill A $=$\linie{25mm}
\fusshilfe{\gfuss{6}{163,2 m²}}
\end{pfaufg}
% DRUCK 2023-OS-K2c deckt 2023-OS-K2c
\begin{pfaufg}{nach~P10\\’23}{7.\textsuperscript{*}}{}
\gtextkaro{\begin{tikzpicture}[line width=0.7pt,font=\small]\draw (0,0) -- (5.16,0) -- (4.08,1.6) -- (1.08,1.6) -- cycle;\draw[densely dotted,line width=0.6pt] (1.08,1.6) -- (1.08,0);\pfrw[0.2]{(1.08,0)}{90}\node[anchor=west,font=\footnotesize] at (1.12,0.85) {$8{,}00$ m};\pfwinkel[0.55][0.6]{(0,0)}{0}{56}{}\node[font=\scriptsize] at (0.74,0.4) {$56^\circ$};\pfwinkel[0.55][0.6]{(5.16,0)}{124}{180}{}\node[font=\scriptsize] at (4.42,0.4) {$56^\circ$};\draw[line width=0.4pt,|-|] (1.08,1.95) -- (4.08,1.95) node[midway,above,inner sep=1pt,font=\footnotesize] {$15{,}00$ m};\draw[line width=0.4pt,|-|] (0,-0.35) -- (5.16,-0.35) node[midway,below,inner sep=1pt,font=\footnotesize] {$25{,}80$ m};\node[font=\scriptsize,color=mbgrau,anchor=north] at (current bounding box.south) {(nicht maßstabsgerecht)};\end{tikzpicture}}{Gegeben ist das folgende Trapez. Berechne den Umfang des Trapezes.}{5}
\par\medskip\noindent\hfill u $=$\linie{25mm}
\fusshilfe{\gfuss{7}{60,1 m}}
\end{pfaufg}
% DRUCK 2014-OS-K5c deckt 2014-OS-K5c
\begin{pfaufg}{nach~P10\\’14}{8.\textsuperscript{*}}{}
\gtextkaro{\begin{tikzpicture}[scale=0.7,line width=0.7pt,font=\small]\fill[black!10] (0,4.4) -- (2.2,4.4) -- (2.2,1.1) -- (1.9,0) -- (0.3,0) -- (0,1.1) -- cycle;\fill[black!22] (0.3,3.6) rectangle (1.9,1.1);\draw (0,4.4) -- (2.2,4.4) -- (1.9,3.6) -- (0.3,3.6) -- cycle;\draw (0,4.4) -- (0,1.1) (2.2,4.4) -- (2.2,1.1);\draw (0.3,3.6) -- (0.3,0) (1.9,3.6) -- (1.9,0);\draw[dashed,line width=0.5pt] (0,1.1) -- (2.2,1.1);\draw (0,1.1) -- (0.3,0) -- (1.9,0) -- (2.2,1.1);\draw[line width=0.5pt] (1.1,4.4) -- (1.1,3.6);\draw[line width=0.4pt] (1.1,3.75) -- (1.25,3.75) -- (1.25,3.6);\draw[line width=0.5pt] (1.1,1.1) -- (1.1,0) node[pos=0.4,left,inner sep=1pt] {$h$};\draw[line width=0.4pt] (1.1,0.15) -- (1.25,0.15) -- (1.25,0);\draw[line width=0.4pt,<->] (0,4.65) -- (2.2,4.65) node[midway,above,inner sep=1pt] {110};\draw[line width=0.4pt,<->] (2.45,4.4) -- (2.45,1.1) node[midway,right,inner sep=1pt] {165};\draw[line width=0.4pt,<->] (2.42,1.02) -- (2.12,-0.08) node[midway,right,inner sep=2pt] {57};\draw[line width=0.4pt,<->] (0.3,-0.25) -- (1.9,-0.25) node[midway,below,inner sep=1pt] {80};\node[font=\scriptsize,color=mbgrau,anchor=north] at (1.1,-0.75) {(Maße in cm; nicht maßstabsgerecht)};\end{tikzpicture}}{Bäckermeister Neumann lässt eine Kühlvitrine in Form eines Prismas bauen (Maße in cm). Als Grundfläche ist ein gleichschenkliges Trapez mit einem Flächeninhalt von 5\,225 cm² vorgesehen. Weise nach, dass die Tiefe $h$ der Vitrine für Kuchenplatten mit einem Durchmesser von 50 cm ausreicht.}{5}
\fusshilfe{\gfuss{8}{ja}}
\end{pfaufg}

% ===================== D Flächeninhalt von Drachen =====================
\gkopfvor{\gbstvor{D}\pfstufekopf{Flächeninhalt von Drachen}{}}
% DRUCK 2025-OS-K2b deckt 2025-OS-K2b
\begin{pfaufg}{nach~P10\\’25}{9.}{}
\gtextkaro{\begin{tikzpicture}[line width=0.7pt,font=\small,scale=0.045]\coordinate (A) at (-25,0);\coordinate (C) at (25,0);\coordinate (D) at (0,25);\coordinate (B) at (0,-55);\coordinate (F) at (0,0);\draw (A) -- (D) -- (C) -- (B) -- cycle;\draw[line width=0.4pt] (A) -- (C) (D) -- (B);\pfrw[0.3]{(F)}{90}\pfrw[0.3]{(D)}{225}\node[anchor=east] at (A) {$A$};\node[anchor=west] at (C) {$C$};\node[anchor=south] at (D) {$D$};\node[anchor=north] at (B) {$B$};\node[anchor=north west] at (F) {$F$};\draw[line width=0.4pt,densely dotted] (-36,25) -- (0,25);\draw[line width=0.4pt,|-|] (-36,0) -- (-36,25);\node[rotate=90,anchor=south,font=\footnotesize] at (-36.5,12.5) {$25$ cm};\end{tikzpicture}}{Die Diagonalen des abgebildeten Drachenvierecks haben die Maße $\overline{AC} = 50$~cm und $\overline{BD} = 80$~cm. Ermittle den Flächeninhalt des Drachenvierecks $ABCD$.}{4}
\par\medskip\noindent\hfill A $=$\linie{25mm}
\fusshilfe{\gfuss{9}{2000 cm²}}
\end{pfaufg}

% ===================== E Kreis =====================
\gkopfvor{\gbstvor{E}\pfstufekopf{Kreis}{}}
% DRUCK 2024-OS-K4a deckt 2024-OS-K4a
\begin{pfaufg}{nach~P10\\’24}{10.}{}
\gtextkaro{\begin{tikzpicture}[line width=0.7pt]\fill[black!12] (0,0) ellipse (0.75 and 0.26);\draw (-0.75,0) -- (0,2.0) -- (0.75,0);\draw (-0.75,0) arc (180:360:0.75 and 0.26);\draw[dashed,line width=0.5pt] (-0.75,0) arc (180:0:0.75 and 0.26);\node[font=\scriptsize,color=mbgrau,anchor=north] at (current bounding box.south) {(nicht maßstabsgerecht)};\end{tikzpicture}}{Eine Firma stellt Kegel her. Der Kegel hat die Maße $r = 30$ cm und $h = 80$ cm. Berechne den Flächeninhalt der Grundfläche des Kegels.}{4}
\par\medskip\noindent\hfill G $=$\linie{25mm}
\fusshilfe{\gfuss{10}{2827,4 cm²}}
\end{pfaufg}
% DRUCK 2016-OS-K3b deckt 2016-OS-K3b
\begin{pfaufg}{nach~P10\\’16}{11.}{}
Der Abwurfring beim Kugelstoßen ist ein Kreis mit einem Durchmesser von ca. 2,14 m. Eine Kugelstoßanlage soll neu gebaut werden. Ermittle dazu den Flächeninhalt des Kreises, den der Abwurfring einschließt.
\karo{4}
\par\medskip\noindent\hfill A $=$\linie{25mm}
\fusshilfe{\gfuss{11}{3,60 m²}}
\end{pfaufg}
% DRUCK 2018-OS-K6a deckt 2018-OS-K6a
\begin{pfaufg}{nach~P10\\’18}{12.}{}
\gtextkaro{\begin{tikzpicture}[line width=0.7pt]\fill[black!6] (0,0) rectangle (4.0,2.4);\draw[densely dotted] (0,0) rectangle (4.0,2.4);\filldraw[fill=black!45,line width=0.8pt] (1.25,1.45) circle (0.6);\node[font=\scriptsize,color=mbgrau,anchor=north] at (current bounding box.south) {(nicht maßstabsgerecht)};\end{tikzpicture}}{Der Fuchsturm ist ein beliebtes Ausflugsziel in der Stadt Jena. Der Turm hat einen Durchmesser $d$ von 6,4 m. Die Rasenfläche um den Turm soll erneuert werden. Sie ist 20,0 m lang und 12,0 m breit. Berechne die zu erneuernde Rasenfläche.}{5}
\par\medskip\noindent\hfill A $=$\linie{25mm}
\fusshilfe{\gfuss{12}{207,8 m²}}
\end{pfaufg}
\end{document}